We study a toy machine-learned potential for 5--13 atom Lennard-Jones clusters (perturbed minima and random configurations) with three input representations (raw coordinates, sorted inverse pair distances, sum-pooled Behler--Parrinello symmetry functions) and two regressors (kernel ridge regression, a small MLP), plus an atomwise Behler--Parrinello-style network, all reimplemented. We measure test energy error, learning curves for 50--1500 training configurations, and error under random rotation and atom permutation of the test set, with 3--5 seeds on CPU. Raw coordinates fail: with randomly rotated and ordered training data, raw KRR is indistinguishable from predicting the mean energy (MAE 12.178 versus 12.203 at $n{=}500$) and the raw MLP is worse (15.028), so their error does not rise further under rotation or permutation; their predictions change by 2.88 and 15.2 energy units on average under such a transform, while every invariant system stays below $10^{-4}$. At $n{=}500$ the sorted-distance MLP is the most accurate system (0.635), followed by symmetry-function KRR (0.889) and sorted-distance KRR (1.095). Symmetry-function KRR is more accurate than sorted-distance KRR at $n{=}50$, 150 and 500 (Welch $p=0.012$, 0.032, 0.026) and less accurate at $n{=}1500$ (0.71 versus 0.45), which supports our pre-registered small-data hypothesis for KRR, on 3 seeds per size at $n{=}50$ and 150. This verdict depends on the KRR hyperparameter grid: a narrower grid, used in the first version of this study, clips the kernel width and gives 1.587 at $n{=}500$. The selected kernels are nearly flat, and plain linear ridge regression on the same descriptor reaches 0.805, whereas the two networks on symmetry functions reach only 4.513 and 3.130. Removing the angular symmetry functions has no resolved effect (0.785 versus 0.889, $p=0.190$). This is a small single-system study with a shared minima library; it supports no claim beyond it.
